% =====================================================================
% Skeleton G — 3 Layer Vertical + Central Hero + N Parallel Instances
% =====================================================================
% USAGE: Adapted from a GOLD STANDARD figure (Federated Traffic Prediction).
% Sub-agent: copy ENTIRE file as figure.tex, change content.
%
% LAYOUT MODEL:
%   ┌────────────────────────────────────────────────┐  ┌───┐
%   │   Top Layer: Central Hero (aggregator/server)  │  │ R │
%   │              [Big purple box]                  │  │ i │
%   ├────────────────────────────────────────────────┤  │ g │
%   │ Mid Layer: N parallel instances (3 cities/    │  │ h │
%   │  workers/clients) each with internal hero      │  │ t │
%   ├────────────────────────────────────────────────┤  │   │
%   │ Bot Layer: Edge/Input + multi-modal sources    │  │ V │
%   │ (camera/sensor/GPS for each instance)          │  │ . │
%   └────────────────────────────────────────────────┘  │ l │
%   Bottom-most: experiment result bar chart            │ . │
%                                                       └───┘
%
% WHEN TO USE: Federated learning, distributed multi-agent systems,
% edge-cloud architectures. Use when you have 1 central + N parallel.
%
% CONTAINS: Central hero (FedAvg server) + 3 parallel cities each with
% ST-GNN (spatial attn + temporal conv) + edge preprocessing per city +
% multi-modal sensors (camera/磁/GPS) + result bar chart + dashed
% feedback rails (上传/下发) + 中文 vertical right label "联邦迭代".
%
% CONSTRAINTS — DO NOT VIOLATE (each one fixed a real bug):
%
%   [RENAMING]
%   - 中心服务器 / 加权 FedAvg 聚合 — RENAME for non-FL topics
%     (e.g., "Coordinator / Multi-Agent Aggregator" for multi-agent LLM).
%   - 3 city names + 3 ST-GNN heroes — RENAME (e.g., 3 agents/workers).
%   - 3 sets of multi-modal sensors (camera/磁/GPS) — RENAME to topic
%     input modalities.
%   - "联邦迭代" right-side vertical label — RENAME to topic loop name.
%   - Bottom bar chart labels AND bar heights — RENAME labels AND
%     recompute bar heights to match the new numeric labels.
%
%   [small_blue BOX TEXT LIMIT — hardcoded 1.6cm width]
%   - small_blue is fixed 1.6cm wide → max ~5 CJK chars OR ~10 English
%     chars. Text like "LLM (GPT-4)" truncates inside the box.
%     Use a single short token ("LLM" or "GPT-4"), not both.
%
%   [ZONE CHIPS (3 layer labels) — position above zones]
%   - All 3 chips use anchor=south west, positioned ABOVE their zone
%     (chip y > zone top y). Do NOT place chips inside the zone's
%     top-left corner — they overlap the content immediately below.
%
%   [CHART PANEL — must enclose ALL its parts]
%   - Panel rectangle MUST span the full y-range of: x-axis arrow tip,
%     AND axis tick labels, AND chart title. Current bounds: y=0.4 to
%     y=4.6 (height 4.2cm); title at y=3.9; "数据不出域" chip at y=3.9.
%   - Bar heights MUST match their numeric labels — recompute if you
%     change data values, do NOT keep old heights.
%
%   [ZONE BOUNDS — leave room for adjacent chips/titles]
%   - Training zone bottom y=8.6 (gap above for 本地训练层 chip).
%   - Edge zone bottom y=4.6 (clearance for chart title at y=3.9).
% =====================================================================

\documentclass[tikz,border=30pt]{standalone}
\usepackage{tikz}
\usepackage[fontset=none]{ctex}
\setCJKmainfont{PingFang SC}
\setCJKsansfont{PingFang SC}
\usepackage{amsmath}
\usetikzlibrary{shapes, arrows.meta, positioning, fit, backgrounds, calc, shadows}

% =====================================================================
% STYLE PRESET: Academic Professional (default)
% =====================================================================
% G's color identity:
%   hero=red (central server), primary=green (parallel instances/cities),
%   accent=orange (sensors), data_curve=blue (chart axes/bars).
%   drawPurple stays HARDCODED — G's "feedback flow" signature color
%   (上传/下发 arrows) which should NOT change across style swaps.
%
% Swap style by replacing PRESET DEFINITION block only.
% See references/style-presets/README.md for token semantics.
% =====================================================================

% ----- BEGIN PRESET DEFINITION -----
\definecolor{c_primary_line}{HTML}{82B366}       \definecolor{c_primary_fill}{HTML}{D5E8D4}
\definecolor{c_hero_line}{HTML}{B85450}          \definecolor{c_hero_fill}{HTML}{F8CECC}
\definecolor{c_accent_line}{HTML}{D79B00}        \definecolor{c_accent_fill}{HTML}{FFE6CC}
\definecolor{c_data_curve_line}{HTML}{6C8EBF}    \definecolor{c_data_curve_fill}{HTML}{DAE8FC}
\definecolor{c_neutral_line}{HTML}{666666}       \definecolor{c_neutral_fill}{HTML}{F5F5F5}
\definecolor{c_bar_1}{HTML}{3B82F6}
% ----- END PRESET DEFINITION -----

% ----- BEGIN BACKWARDS-COMPAT ALIASES (DO NOT EDIT) -----
\colorlet{drawBlueLine}{c_data_curve_line}       \colorlet{drawBlueFill}{c_data_curve_fill}
\colorlet{drawGreenLine}{c_primary_line}         \colorlet{drawGreenFill}{c_primary_fill}
\colorlet{drawOrangeLine}{c_accent_line}         \colorlet{drawOrangeFill}{c_accent_fill}
\colorlet{drawRedLine}{c_hero_line}              \colorlet{drawRedFill}{c_hero_fill}
\colorlet{drawGreyLine}{c_neutral_line}          \colorlet{drawGreyFill}{c_neutral_fill}
\colorlet{barBlue}{c_bar_1}
% ----- END BACKWARDS-COMPAT ALIASES -----

% G-internal signature color (NOT preset-controlled): the "上传/下发"
% federated feedback flow stays purple across all styles for visual
% consistency of the FL loop semantics.
\definecolor{drawPurpleLine}{HTML}{9673A6}       \definecolor{drawPurpleFill}{HTML}{E1D5E7}

\pgfdeclarelayer{background}
\pgfsetlayers{background,main}

\begin{document}
\begin{tikzpicture}[
    every node/.style={font=\small},
    base_box/.style={rectangle, rounded corners=3pt, align=center,
        minimum height=0.9cm, minimum width=2.8cm,
        inner sep=10pt, drop shadow={opacity=0.15}, thick},
    green_node/.style={base_box, fill=drawGreenFill, draw=drawGreenLine},
    red_node/.style={base_box, fill=drawRedFill, draw=drawRedLine, font=\small\bfseries},
    grey_node/.style={base_box, fill=drawGreyFill, draw=drawGreyLine},
    small_blue/.style={rectangle, rounded corners=2pt, align=center,
        minimum height=0.65cm, minimum width=1.6cm,
        inner sep=5pt, thick, font=\footnotesize,
        fill=drawBlueFill!60, draw=drawBlueLine!80},
    sensor/.style={rectangle, rounded corners=2pt, align=center,
        minimum height=0.5cm, minimum width=1.2cm,
        inner sep=3pt, thick, font=\scriptsize,
        fill=drawOrangeFill, draw=drawOrangeLine},
    arrow/.style={-{Stealth[scale=0.9]}, thick, color=black!70,
        shorten >=2pt, shorten <=1pt},
    thin_arrow/.style={-{Stealth[scale=0.7]}, thick, color=drawGreyLine},
]

% ================================================================
%  LAYER 3 — FEDERATED AGGREGATION (top)
% ================================================================
\node[red_node, minimum width=7.5cm, minimum height=2.4cm] (server) at (6, 16.5)
    {\textbf{中心服务器}\\[2pt]{\footnotesize 加权 FedAvg 聚合}\\[-1pt]{\footnotesize 生成并下发全局模型}};

% ================================================================
%  LAYER 2 — LOCAL TRAINING (middle)
% ================================================================
% --- City A ---
\node[green_node, minimum width=4.4cm, minimum height=3.2cm] (localA) at (0, 10.5) {};
\node[font=\small\bfseries, text=drawGreenLine!80!black] at (0, 11.6) {城市A\,$\cdot$\,ST-GNN};
\node[small_blue] (spatA) at (-0.8, 10.2) {空间注意力};
\node[small_blue] (tempA) at (0.8, 10.2) {时间卷积};
\draw[thin_arrow] (spatA.east) -- (tempA.west);
\node[font=\scriptsize, text=drawGreyLine, align=center] at (0, 9.3)
    {路网拓扑\,$\to$\,周期模式};

% --- City B ---
\node[green_node, minimum width=4.4cm, minimum height=3.2cm] (localB) at (6, 10.5) {};
\node[font=\small\bfseries, text=drawGreenLine!80!black] at (6, 11.6) {城市B\,$\cdot$\,ST-GNN};
\node[small_blue] (spatB) at (5.2, 10.2) {空间注意力};
\node[small_blue] (tempB) at (6.8, 10.2) {时间卷积};
\draw[thin_arrow] (spatB.east) -- (tempB.west);
\node[font=\scriptsize, text=drawGreyLine, align=center] at (6, 9.3)
    {路网拓扑\,$\to$\,周期模式};

% --- City C ---
\node[green_node, minimum width=4.4cm, minimum height=3.2cm] (localC) at (12, 10.5) {};
\node[font=\small\bfseries, text=drawGreenLine!80!black] at (12, 11.6) {城市C\,$\cdot$\,ST-GNN};
\node[small_blue] (spatC) at (11.2, 10.2) {空间注意力};
\node[small_blue] (tempC) at (12.8, 10.2) {时间卷积};
\draw[thin_arrow] (spatC.east) -- (tempC.west);
\node[font=\scriptsize, text=drawGreyLine, align=center] at (12, 9.3)
    {路网拓扑\,$\to$\,周期模式};

% ================================================================
%  LAYER 1 — EDGE COLLECTION (bottom)
% ================================================================
% Preprocessing
\node[grey_node, minimum width=3.2cm] (prepA) at (0, 7.0)
    {\footnotesize 本地预处理};
\node[grey_node, minimum width=3.2cm] (prepB) at (6, 7.0)
    {\footnotesize 本地预处理};
\node[grey_node, minimum width=3.2cm] (prepC) at (12, 7.0)
    {\footnotesize 本地预处理};

% Sensors — City A
\node[sensor] (camA) at (-1.5, 5.1) {摄像头};
\node[sensor] (magA) at (0, 5.1) {地磁};
\node[sensor] (gpsA) at (1.5, 5.1) {GPS};
% Sensors — City B
\node[sensor] (camB) at (4.5, 5.1) {摄像头};
\node[sensor] (magB) at (6, 5.1) {地磁};
\node[sensor] (gpsB) at (7.5, 5.1) {GPS};
% Sensors — City C
\node[sensor] (camC) at (10.5, 5.1) {摄像头};
\node[sensor] (magC) at (12, 5.1) {地磁};
\node[sensor] (gpsC) at (13.5, 5.1) {GPS};

% ================================================================
%  EXPERIMENT RESULTS PANEL (bottom)
% ================================================================
% Panel resized: height 3.8 → 4.2 (center stays 1.7) so the panel spans
% y=-0.4..3.8 (was -0.2..3.6). Previous panel top y=3.6 was EXACTLY at
% the y-axis arrow tip (scope y=3.4 + shift y=0.20 = abs y=3.6) — the
% arrowhead touched the panel border. New panel top y=3.8 gives 0.2cm
% clearance. Title moved from y=3.7 → 3.9 to stay ABOVE the new panel
% top (3.8).
\node[rectangle, rounded corners=6pt, draw=drawBlueLine, fill=drawBlueFill!15,
    minimum width=15cm, minimum height=4.2cm, thick, inner sep=0pt] (panel) at (6, 1.7) {};
\node[font=\small\bfseries, text=drawBlueLine!80!black, anchor=south west]
    at (-0.8, 3.9) {实验结果：三城市联合预测性能};

% Embedded bar chart
% Fix G.6: was y-axis labeled 5/10/15/20 (max 20) but bar value labels
% said 22.1/25.3 — bar heights and axis scale didn't match (visual lie).
% Now: extend axis to label "30", scale = 0.1 cm per 1 unit of value,
% bar heights = 0.2 + 0.1 * value (correct linear mapping).
\begin{scope}[shift={(1.5, 0.20)}]
  % Axes — extended to y=3.2 (label 30)
  \draw[-{Stealth[scale=0.6]}, thick, drawGreyLine] (0.5, 0.2) -- (0.5, 3.4);
  \draw[-{Stealth[scale=0.6]}, thick, drawGreyLine] (0.5, 0.2) -- (5.5, 0.2);
  % Grid lines + Y-axis ticks at 5/10/15/20/25/30
  \foreach \y/\lab in {0.7/5, 1.2/10, 1.7/15, 2.2/20, 2.7/25, 3.2/30} {
    \draw[gray!20, thin] (0.5, \y) -- (5.2, \y);
    \node[font=\scriptsize, left, text=drawGreyLine] at (0.4, \y) {\lab};
  }

  % MAE group — bar heights NOW match axis (22.1→y=2.41, 18.1→y=2.01)
  \fill[drawGreyFill, draw=drawGreyLine] (1.0, 0.2) rectangle (1.6, 2.41);
  \fill[barBlue, draw=barBlue!80!black, rounded corners=1pt] (1.8, 0.2) rectangle (2.4, 2.01);
  \node[font=\scriptsize, text=drawGreyLine] at (1.3, 2.56) {22.1};
  \node[font=\scriptsize, text=barBlue!80!black] at (2.1, 2.16) {18.1};
  \node[font=\scriptsize] at (1.7, -0.05) {MAE};

  % RMSE group — 25.3→y=2.73, 19.8→y=2.18
  \fill[drawGreyFill, draw=drawGreyLine] (3.0, 0.2) rectangle (3.6, 2.73);
  \fill[barBlue, draw=barBlue!80!black, rounded corners=1pt] (3.8, 0.2) rectangle (4.4, 2.18);
  \node[font=\scriptsize, text=drawGreyLine] at (3.3, 2.88) {25.3};
  \node[font=\scriptsize, text=barBlue!80!black] at (4.1, 2.33) {19.8};
  \node[font=\scriptsize] at (3.7, -0.05) {RMSE};

  % Improvement arrows + labels — to the RIGHT of each pair
  \draw[-{Stealth[scale=0.5]}, very thick, drawRedLine] (2.55, 2.41) -- (2.55, 2.01);
  \node[font=\scriptsize\bfseries, text=drawRedLine, anchor=west]
      at (2.6, 2.21) {$\downarrow$18.3\%};
  \draw[-{Stealth[scale=0.5]}, very thick, drawRedLine] (4.55, 2.73) -- (4.55, 2.18);
  \node[font=\scriptsize\bfseries, text=drawRedLine, anchor=west]
      at (4.6, 2.46) {$\downarrow$21.7\%};
\end{scope}

% Legend (right side of panel)
\begin{scope}[shift={(8.5, 1.5)}]
  \fill[drawGreyFill, draw=drawGreyLine] (0, 0.7) rectangle (0.4, 0.9);
  \node[font=\scriptsize, anchor=west] at (0.5, 0.8) {基线方法};
  \fill[barBlue, draw=barBlue!80!black] (0, 0.25) rectangle (0.4, 0.45);
  \node[font=\scriptsize, anchor=west] at (0.5, 0.35) {FedTraffic};
\end{scope}

% Key insight — Fix G.5: was green-on-green chip inside blue panel
% (color conflict + visually a "data point" inside the chart area).
% Move OUTSIDE the panel (above it, anchor=south) and harmonize color
% with the panel's blue theme so it reads as a summary annotation.
% Moved 3.7 → 3.9 to stay above the resized panel top (now y=3.8) — same
% as the title repositioning.
\node[font=\footnotesize\bfseries, text=drawBlueLine!80!black,
      fill=drawBlueFill!40, draw=drawBlueLine!50, rounded corners=3pt,
      inner sep=5pt, anchor=south west] at (9.0, 3.9)
    {数据不出域，隐私安全};

% ================================================================
%  CONNECTIONS — Sensors → Preprocessing
% ================================================================
\draw[arrow] (camA.north) -- ([xshift=-10pt]prepA.south);
\draw[arrow] (magA.north) -- (prepA.south);
\draw[arrow] (gpsA.north) -- ([xshift=10pt]prepA.south);

\draw[arrow] (camB.north) -- ([xshift=-10pt]prepB.south);
\draw[arrow] (magB.north) -- (prepB.south);
\draw[arrow] (gpsB.north) -- ([xshift=10pt]prepB.south);

\draw[arrow] (camC.north) -- ([xshift=-10pt]prepC.south);
\draw[arrow] (magC.north) -- (prepC.south);
\draw[arrow] (gpsC.north) -- ([xshift=10pt]prepC.south);

% ================================================================
%  CONNECTIONS — Preprocessing → Local Models
% ================================================================
\draw[arrow] (prepA.north) -- (localA.south);
\draw[arrow] (prepB.north) -- (localB.south);
\draw[arrow] (prepC.north) -- (localC.south);

% ================================================================
%  UPLOAD GRADIENTS (orange thick, L-shaped convergence)
% ================================================================
% A → server: up, right, up
\draw[-{Stealth[scale=1.1]}, very thick, color=drawOrangeLine, rounded corners=6pt]
    ([xshift=-10pt]localA.north) -- ++(0, 1.5) -| (4.0, 15.3);
% B → server: straight up
\draw[-{Stealth[scale=1.1]}, very thick, color=drawOrangeLine]
    ([xshift=-10pt]localB.north) -- ([xshift=-10pt]server.south);
% C → server: up, left, up
\draw[-{Stealth[scale=1.1]}, very thick, color=drawOrangeLine, rounded corners=6pt]
    ([xshift=-10pt]localC.north) -- ++(0, 1.5) -| (8.0, 15.3);

% Upload label
\node[font=\footnotesize\bfseries, text=drawOrangeLine,
      fill=white, fill opacity=0.9, text opacity=1,
      inner sep=3pt, rounded corners=2pt] at (1.5, 13.7) {上传梯度};

% ================================================================
%  DOWNLOAD MODEL (blue dashed, L-shaped divergence)
% ================================================================
% server → A: down, left, down
\draw[-{Stealth[scale=0.9]}, thick, dashed, color=drawBlueLine, rounded corners=6pt]
    (4.5, 15.3) -- ++(0, -1.3) -| ([xshift=10pt]localA.north);
% server → B: straight down
\draw[-{Stealth[scale=0.9]}, thick, dashed, color=drawBlueLine]
    ([xshift=10pt]server.south) -- ([xshift=10pt]localB.north);
% server → C: down, right, down
\draw[-{Stealth[scale=0.9]}, thick, dashed, color=drawBlueLine, rounded corners=6pt]
    (7.5, 15.3) -- ++(0, -1.3) -| ([xshift=10pt]localC.north);

% Download label
\node[font=\footnotesize\bfseries, text=drawBlueLine,
      fill=white, fill opacity=0.9, text opacity=1,
      inner sep=3pt, rounded corners=2pt] at (10.5, 13.7) {下发全局模型};

% ================================================================
%  ZONE BACKGROUNDS (pgfonlayer)
% ================================================================
\begin{pgfonlayer}{background}
  % Aggregation zone
  \fill[drawRedFill!20, rounded corners=10pt] (-3.2, 14.7) rectangle (15.2, 18.3);
  \draw[drawRedLine!50, dashed, thick, rounded corners=10pt] (-3.2, 14.7) rectangle (15.2, 18.3);
  % Training zone — bottom extended 8.3 → 8.6 so the 边缘采集层 chip
  % above Edge zone (chip top at y=8.6) doesn't poke into Training zone.
  \fill[drawGreenFill!20, rounded corners=10pt] (-3.2, 8.6) rectangle (15.2, 12.7);
  \draw[drawGreenLine!50, dashed, thick, rounded corners=10pt] (-3.2, 8.6) rectangle (15.2, 12.7);
  % Edge collection zone — bottom raised 4.3 → 4.6 to give the
  % "实验结果" chart-panel title more vertical breathing room (was 0.3cm
  % gap from title top y≈4.0 to zone dashed border at 4.3, now 0.6cm).
  \fill[drawOrangeFill!20, rounded corners=10pt] (-3.2, 4.6) rectangle (15.2, 8.0);
  \draw[drawOrangeLine!50, dashed, thick, rounded corners=10pt] (-3.2, 4.6) rectangle (15.2, 8.0);
\end{pgfonlayer}

% ================================================================
%  ZONE LABELS (capsule style — all 3 chips OUTSIDE their zones above)
% ================================================================
% All three chips share the same style: anchor=south west placed JUST
% ABOVE each zone's top border, with chip bottom touching zone top.
% Previously 联邦聚合层 was anchor=north west INSIDE its zone at top-left,
% while 本地训练层 / 边缘采集层 were anchor=south west ABOVE their zones
% — inconsistent visual treatment. Now uniformly above-zone.
% (本地训练层/边缘采集层 moved above to avoid overlap with City A / Prep
% box content; Training zone bottom extended 8.3→8.6 to fit 边缘采集层
% chip in the inter-zone gap.)
\node[font=\footnotesize\bfseries, text=drawRedLine!80!black,
      fill=drawRedFill!60, draw=drawRedLine!40, rounded corners=3pt, inner sep=4pt,
      anchor=south west] at (-2.7, 18.35) {联邦聚合层};
\node[font=\footnotesize\bfseries, text=drawGreenLine!80!black,
      fill=drawGreenFill!60, draw=drawGreenLine!40, rounded corners=3pt, inner sep=4pt,
      anchor=south west] at (-2.7, 12.75) {本地训练层};
\node[font=\footnotesize\bfseries, text=drawOrangeLine!80!black,
      fill=drawOrangeFill!60, draw=drawOrangeLine!40, rounded corners=3pt, inner sep=3pt,
      anchor=south west] at (-2.7, 8.05) {边缘采集层};

% ================================================================
%  DECORATIVE — Iteration cycle annotation (right margin)
% ================================================================
\draw[-{Stealth[scale=0.8]}, thick, color=drawPurpleLine, rounded corners=4pt]
    (15.3, 12.5) -- (15.3, 17.5);
\draw[-{Stealth[scale=0.8]}, thick, dashed, color=drawPurpleLine, rounded corners=4pt]
    (15.7, 17.5) -- (15.7, 12.5);
\node[font=\scriptsize, text=drawPurpleLine, anchor=west] at (16.0, 15.0)
    {\shortstack{联\\邦\\迭\\代}};

\end{tikzpicture}
\end{document}
